% Partie : logique propositionnelle
% Chapitre : syntaxe et sémantique
% Section : complétude

\documentclass[../../../include/open-logic-section]{subfiles}

\begin{document}

\olfileid[fr]{pl}{prp}{com}

\olsection{Complétude de la logique propositionnelle}

Nous considérons ici le système axiomatique classique
pour les connecteurs propositionnels, avec le modus ponens
comme règle d'inférence. En particulier, le théorème de
la déduction et l'élimination de la double négation y sont
disponibles. Une preuve utilise toujours un nombre fini
d'hypothèses, même si l'ensemble d'hypothèses est infini.

\begin{defn}
\ollabel{def:MaxCon}
Un ensemble~$\Gamma$ de !!{formula}s est
\emph{maximalement cohérent} s'il est cohérent
et si tout ensemble cohérent~$\Delta$ tel que
$\Gamma\subseteq\Delta$ est égal à~$\Gamma$.
\end{defn}

Nous utiliserons le fait d'extension suivant :
si $\Gamma$ est cohérent, alors, pour toute formule~$!A$,
au moins l'un des deux ensembles
$\Gamma\cup\{!A\}$ et $\Gamma\cup\{\lnot !A\}$
est cohérent. En effet, si tous deux étaient incohérents,
le théorème de la déduction donnerait, à partir de~$\Gamma$,
$!A\lif\lfalse$ et
$\lnot !A\lif\lfalse$. Les axiomes classiques
permettraient alors de déduire, à partir de~$\Gamma$,
$\lnot !A$ et $\lnot\lnot !A$, donc aussi~$!A$.
On obtiendrait ainsi une contradiction à partir de~$\Gamma$,
alors que cet ensemble est cohérent.

\begin{lem}[Lemme de vérité]
  \ollabel{lem:truth}
Soit $\Gamma$ maximalement cohérent. Alors :
\begin{enumerate}
\item \ollabel{lem:truth-1} $!A \in \Gamma$ si et seulement si
  $\lnot !A \notin \Gamma$ ;
\item \ollabel{lem:truth-2} $\Gamma \Proves !A$ si et seulement si
  $!A \in \Gamma$ ;
\item \ollabel{lem:truth-3} $!A \lif !B \in \Gamma$ si et seulement si
  $!A \notin \Gamma$ ou $!B \in \Gamma$.
\end{enumerate}
\end{lem}

\begin{proof}
\begin{enumerate}
\item Si $!A\in\Gamma$, la cohérence exclut
  $\lnot !A\in\Gamma$. Si $!A$ et $\lnot !A$
  étaient toutes deux absentes de~$\Gamma$, le fait
  d'extension fournirait un sur-ensemble cohérent
  strict de~$\Gamma$. Cela contredirait sa maximalité.
\item Si $!A\in\Gamma$, l'axiome d'hypothèse donne
  $\Gamma\Proves !A$. Si $\Gamma\Proves !A$
  mais $!A\notin\Gamma$, le point~\olref{lem:truth-1}
  donne $\lnot !A\in\Gamma$. On aurait ainsi
  $\Gamma\Proves !A$ et
  $\Gamma\Proves\lnot !A$, ce qui contredit
  la cohérence.
\item Pour le sens direct, appliquer le modus ponens
  si $!A\in\Gamma$. Pour la réciproque, distinguer
  les cas $!B\in\Gamma$ et $!A\notin\Gamma$ ;
  dans le second, la première assertion fournit
  $\lnot !A\in\Gamma$. Les axiomes classiques
  donnent respectivement
  $!B\lif(!A\lif !B)$ et
  $\lnot !A\lif(!A\lif !B)$.
  Les détails des deux applications du modus ponens
  sont laissés en exercice.
\end{enumerate}
\end{proof}

Le sens de droite à gauche du
\olref{lem:truth}\olref{lem:truth-2}
est l'axiome d'hypothèse ; le sens de gauche à droite
montre que $\Gamma$ est \emph{clos déductivement} :
si $\Gamma\Proves !A$, alors $!A\in\Gamma$.

\begin{prob}
  Détailler le dernier cas de la preuve du
  \olref{lem:truth}, puis vérifier directement
  les clauses correspondant à la conjonction et
  à la disjonction dans le système axiomatique choisi.
\end{prob}

\begin{thm}[Complétude]
\ollabel{thm:completeness}
Si $\Gamma$ est cohérent, alors $\Gamma$ est satisfaisable.
\end{thm}

\begin{proof}
Comme les variables propositionnelles sont dénombrables
et les formules sont finies, énumérons toutes les
!!{formula}s du langage sous la forme
$!A_0$, $!A_1$,~\dots.
Définissons par récurrence une suite croissante
d'ensembles de !!{formula}s
$\Gamma_0$, $\Gamma_1$,~\dots :
\begin{align*}
\Gamma_0 & = \Gamma\\
\Gamma_{n+1} & =
\begin{cases}
  \Gamma_n \cup \{!A_n\}
    & \text{si $\Gamma_n\cup\{!A_n\}$ est cohérent;}\\
  \Gamma_n \cup \{\lnot !A_n\}
    & \text{sinon.}
\end{cases}
\end{align*}
Posons ensuite
\[
\Gamma^*=\bigcup_{0\le n}\Gamma_n.
\]
Trois faits achèvent la construction :
\begin{enumerate}
\item Chaque $\Gamma_n$ est cohérent, par récurrence
  sur~$n$ et par le fait d'extension démontré ci-dessus.
\item $\Gamma^*$ est cohérent. En effet, une preuve de
  contradiction n'emploierait qu'un nombre fini
  de formules de~$\Gamma^*$ ; toutes figureraient
  dans un même $\Gamma_n$, ce qui contredirait
  le premier fait.
\item $\Gamma^*$ est maximalement cohérent.
  Pour chaque~$!A_n$, la construction contient
  $!A_n$ ou $\lnot !A_n$. Un sur-ensemble
  cohérent strict devrait ajouter une formule
  absente de~$\Gamma^*$, alors que sa négation
  y figure déjà.
\end{enumerate}
Définissons une valuation~$v$ par
$v(p_i)=\True$ si et seulement si
$p_i\in\Gamma^*$. Par récurrence sur la
structure de~$!A$, le lemme de vérité donne
$\overline v(!A)=\True$ si et seulement si
$!A\in\Gamma^*$. Pour la négation, on utilise
la première assertion du lemme ; pour l'implication,
la troisième. Les clauses de conjonction et de
disjonction découlent des axiomes correspondants
  et de la clôture déductive : une conjonction appartient
à~$\Gamma^*$ exactement quand ses deux conjoints
y appartiennent, et une disjonction exactement
quand l'un des deux disjoints y appartient.
Les constantes logiques se traitent par leurs
axiomes ; les connecteurs définis suivent leurs
abréviations. Ainsi $v$ satisfait~$\Gamma^*$.
Comme $\Gamma\subseteq\Gamma^*$, il satisfait
aussi~$\Gamma$.
\end{proof}

\begin{cor}
Si $\Gamma\Entails !A$, alors $\Gamma\Proves !A$.
\end{cor}

\begin{proof}
Raisonnons par contraposée. Si
$\Gamma\Proves/ !A$, alors
$\Gamma\cup\{\lnot !A\}$ est cohérent :
son incohérence, par le théorème de la déduction
et l'axiome de double négation, entraînerait
$\Gamma\Proves !A$. Par le théorème de complétude,
$\Gamma\cup\{\lnot !A\}$ possède une valuation
satisfaisante. Celle-ci satisfait~$\Gamma$ et
rend~$!A$ fausse ; donc
$\Gamma\Entails/ !A$.
\end{proof}

\begin{prop}[Théorème de compacité]
\ollabel{prop:compactness}
$\Gamma$ est satisfaisable si et seulement si
tout sous-ensemble \emph{fini}~$\Gamma_0$
de~$\Gamma$ est satisfaisable.
\end{prop}

\begin{proof}
Le théorème de correction et celui de complétude
donnent : $\Gamma$ est insatisfaisable si et
seulement s'il est incohérent. Une dérivation
d'incohérence ne fait intervenir qu'un nombre
fini d'hypothèses ; il existe donc un sous-ensemble
fini~$\Gamma_0$ déjà incohérent, et par les mêmes
théorèmes déjà insatisfaisable. La réciproque
est immédiate, car une valuation satisfaisant
$\Gamma$ satisferait chacun de ses sous-ensembles.
\end{proof}

\begin{cor}
$\Gamma\Entails !A$ si et seulement s'il existe
un sous-ensemble fini~$\Gamma_0$ de~$\Gamma$
tel que $\Gamma_0\Entails !A$.
\end{cor}

\end{document}
